\begin{itemize}
\item \textbf{H-A (pore-shape knockdown) -- rejected in its original form.} Circular and square coarse pores at equal net section are equally strong (16.1 vs 17.0\,N/m; predicted 11.3). The net-section factor survived, the shape factor did not.
\item \textbf{H-C (load-path equivalence controls the mode) -- rejected.} Dispersing ligament widths by 15\% lowered strength in both coarse meshes and left the single-avalanche mode unchanged (square) or shortened the tail (round). Replaced by the fibre-bundle load-transfer criterion.
\item \textbf{H-D (disorder optimum for energy absorption) -- not supported after replication.} The Voronoi ``optimum'' at 60\% regularity was a single-seed outlier (14.4 vs 9.2/9.9\,N/m); jittered meshes lose strength without a reproducible gain in work ($1.07\pm0.31$ vs $1.16\pm0.30$\,J/m$^2$ for the ordered lattice across registries).
\item \textbf{H-F$'$ (veins as crack arrestors) -- rejected.} Load-perpendicular veins neither arrest cracks nor add strength (12.0\,N/m, equal to the fine mesh); cracks channel along the columns next to them.
\item \textbf{Flaw halo (P11) -- no shielding.} Pore rings around a hole cost strength relative to the bare hole (16.2 vs 19.1\,N/m); they only reduce the penalty relative to uniformly distributed porosity.
\item \textbf{Gradients (P10) -- detrimental.} Pore-size gradients along or radial to the load are weaker than the uniform mesh (10.1/10.0 vs 11.7\,N/m) because the largest pores define the weakest section.
\item \textbf{Node connectivity (P9) -- minor.} Six-connected triangulated struts are not stronger than four-connected ones; a continuous strut family along the load is what matters (11.5 vs 7.4 vs 6.0\,N/m).
\item \textbf{Unstable and strongly reconstructed candidates.} Every generated geometry relaxed to a mechanically continuous structure, but 6\,\AA{} struts, 20\,\AA{}-vein meshes with residual 2-atom pores, and the 0.40-porosity mesh reconstructed by more than 1\,\AA{} (flagged ``strongly reconstructed''); their results are reported with that flag. No structure fragmented, and none is claimed to be synthesizable.
\item \textbf{Counter-examples worth keeping.} The 45$^\circ$ slit lattice sheds load at 17\% strain and then re-stiffens to $>15$\,N/m at 30\% by ligament rotation (geometric compliance); the 30\,\AA{} crack is not weaker than the 20\,\AA{} crack; the fine period-12 mesh is stronger than the period-16 mesh at equal porosity, but the difference (13.3 vs 11.7\,N/m) is within the $\approx$12\% registry effect measured in Stage~6.
\end{itemize}
